% Diagram of the proof of the existence of regular families of twists (Part 4, Sections cg-selmer, cg-first, cg-lifting).
\begin{figure}[tbp]
\centering
\begin{tikzpicture}[x=1cm,y=1cm]
% row 0: the conclusion
\node[pfnode, text width=3.0cm] (reg) at (-1.4,0) {Thm.~\ref{cg-selmer:thm:regular}, Prop.~\ref{cg-selmer:prop:oneprime}\\regular families};
% row 1
\node[pfnode] (twist) at (-5.2,-1.7) {Thm.~\ref{cg-selmer:thm:twisting}\\twisting theorem};
\node[pfnode] (denom) at (2.6,-1.7) {Thm.~\ref{cg-first:thm:main}\\bound for\\denominators};
% row 2
\node[pfnode, text width=3.4cm] (general) at (-5.2,-3.4) {Thm.~\ref{cg-selmer:thm:general}, Prop.~\ref{cg-selmer:prop:line}\\general twisting theorem};
\node[pfnode] (rows) at (-1.4,-3.4) {Prop.~\ref{cg-selmer:prop:rows}\\main summands};
\node[pfnode, text width=2.6cm] (oper) at (2.6,-3.4) {Props.~\ref{cg-first:prop:operator}, \ref{cg-first:prop:one}\\operator from\\one congruence};
% row 3
\node[pfnode] (descent) at (-6.3,-5.1) {Thm.~\ref{cg-selmer:thm:descent}\\descent};
\node[pfnode] (roots) at (-3.6,-5.1) {Thm.~\ref{cg-selmer:thm:roots}\\roots of unity};
\node[pfnode, text width=2.6cm] (places) at (0.2,-5.1) {Lem.~\ref{cg-first:lem:chain},\\Prop.~\ref{cg-first:prop:places}\\Taylor--Wiles places};
\node[pfnode, text width=2.6cm] (tie) at (3.3,-5.1) {Thms.~\ref{cg-first:thm:tie}, \ref{cg-first:thm:top}\\the tie;\\top coefficient};
\node[pfnode] (ribet) at (6.3,-5.1) {Thm.~\ref{cg-first:thm:ribet}\\Ribet's element};
% row 4
\node[pffact, text width=1.9cm] (cheb) at (-4.9,-6.8) {Fact~\ref{cg-selmer:fact:cheb}\\Chebotarev};
\node[pffact, text width=2.1cm] (galois) at (-2.1,-6.8) {Fact~\ref{cg-selmer:fact:galois}\\Galois\\representations};
\node[pffact, text width=2.9cm] (chaintw) at (1.0,-6.8) {Facts~\ref{cg-first:fact:chain}, \ref{cg-first:fact:tw}\\spherical chains;\\Taylor--Wiles freeness};
\node[pffact, text width=2.4cm] (oddfit) at (4.3,-6.8) {Facts~\ref{cg-first:fact:odd}, \ref{cg-first:fact:fitting}\\odd classes;\\Fitting bound;\\Brauer--Nesbitt};
% row 5
\node[pffact] (gw) at (-6.3,-8.5) {Fact~\ref{cg-selmer:fact:gw}\\Greenberg--Wiles};
\node[pfnode, text width=2.9cm] (loss) at (3.0,-8.5) {Thm.~\ref{cg-lifting:thm:main}, Prop.~\ref{cg-lifting:prop:tatedual}\\local loss above~$\ell$};
% row 6
\node[pffact, text width=1.9cm] (duality) at (0.3,-10.2) {Fact~\ref{cg-lifting:fact:duality}\\local duality};
\node[pfnode, text width=1.9cm] (cuts) at (3.0,-10.2) {Thm.~\ref{cg-lifting:thm:cuts}\\cuts};
\node[pfnode, text width=2.1cm] (generic) at (5.7,-10.2) {\S\ref{cg-lifting:sec:generic}\\generic primes,\\tame types};
% row 7
\node[pfnode, text width=1.9cm] (walk) at (4.35,-11.9) {\S\ref{cg-lifting:sec:walk}\\safe moves,\\the walk};
% row 8
\node[pffact] (frame) at (4.35,-13.6) {Fact~\ref{cg-lifting:fact:frame}\\Kummer charts,\\Kisin modules};
% arrows: premise -> conclusion
\draw[pfedge] (twist) -- (reg);
\draw[pfedge] (rows) -- (reg);
\draw[pfedge] (denom) -- (reg);
\draw[pfedge] (general) -- (twist);
\draw[pfedge] (descent) -- (general);
\draw[pfedge] (roots) -- (general);
\draw[pfedge] (cheb) -- (descent);
\draw[pfedge] (gw) -- (descent);
\draw[pfedge] (galois) -- (rows);
\draw[pfedge] (oper) -- (denom);
\draw[pfedge] (places) -- (oper);
\draw[pfedge] (tie) -- (oper);
\draw[pfedge] (ribet) -- (oper);
\draw[pfedge] (galois) -- (places);
\draw[pfedge] (chaintw) -- (places);
\draw[pfedge] (chaintw) -- (tie);
\draw[pfedge] (loss.east) -| (ribet.south);
\draw[pfedge] (oddfit) -- (ribet);
\draw[pfedge] (duality) -- (loss);
\draw[pfedge] (cuts) -- (loss);
\draw[pfedge] (generic) -- (loss);
\draw[pfedge] (walk) -- (cuts);
\draw[pfedge] (walk) -- (generic);
\draw[pfedge] (frame) -- (walk);
\end{tikzpicture}
\caption{The existence of regular families of twists (Part~4,
Sections~\ref{sec:cg-selmer}--\ref{sec:cg-lifting}), used in the proof of Theorem~C
(Figure~\ref{fig:part4a}). Each box is a result, with its number and a short name; a box with
several numbers, or with the number of a subsection, stands for a group of results. An arrow from
one box to another means that the proof of the second, as written in this paper, uses the first.
Dashed grey boxes are background results admitted without proof.}
\label{fig:part4b}
\end{figure}
